
\def\PUDcodigo{ACE006}

\if\COMPILAR0
    \def\PUDcodacad{04507.7}
\else
    \def\PUDcodacad{04506.7}
\fi

\def\PUDdisciplina{Cálculo 2}

\def\PUDsemestre{S2}

\def\PUDhoras{80}

%NOVOS ---------------------------------------------------
\def\PUDhorasTeoria{80}
\def\PUDhorasPratica{0}

\def\PUDinterECA{---}

\def\PUDatuaisECA{---}

\def\PUDrecursos{Projetor multimídia; Quadro branco e pincel.}

%----------------------------------------------------------

\def\PUDcreditos{4}

%\def\PUDprerequisito{ \hyperref[PUD:ACE002]{ACE002}  \hyperref[PUD:ACE003]{ACE003} }
\def\PUDprerequisito{ \hyperref[PUD:ACE003]{ACE003} }

%\def\PUDpreacad{ \hyperref[PUD:ACE002]{04507.1}  \hyperref[PUD:ACE003]{04507.2} }
\if\COMPILAR0
    \def\PUDpreacad{ \hyperref[PUD:ACE003]{Cálculo 1 (04507.2)} }
\else
    \def\PUDpreacad{ \hyperref[PUD:ACE003]{Cálculo 1 (04506.2)} }
\fi

\def\PUDobjetivos{Compreender o uso da integral na Engenharia através dos seguintes objetivos específicos: Calcular a área de uma região no plano, o volume de um sólido de revolução, o  comprimento de arco de uma curva plana e área de uma superfície de revolução;  Definir a função logarítmica natural, a função exponencial, e as funções  trigonométricas inversas, determinando a derivada e a integral das mesmas;  Definir as funções hiperbólicas, calculando suas derivadas;  Determinar as funções primitivas pelas técnicas de integração;  Calcular limites indeterminados e integrais impróprias;  Representar um ponto e curvas num sistema de coordenadas polares, esboçando gráficos de curvas.}

\def\PUDementa{Aplicações da integral definida.  Funções transcendentes.  Funções hiperbólicas.  As técnicas de integração.  Integrais impróprias.  Noções de coordenadas polares.}

\def\PUDconteudo{ &  Aplicação da integral definida: Área entre curvas;  Volume de sólidos de revolução;  Comprimento de arco de curvas;  Área de superfície de revolução. 
\\ 
 &  Funções Transcendentes: A função logarítmica natural;  A integral e a derivada da função logarítmica natural;  A função exponencial natural;  A derivada e a integral da exponencial;  As funções logarítmica e exponencial num base qualquer;  As funções trigonométricas inversas;  Derivadas das funções trigonométricas inversas. 
\\ 
 &  Funções Hiperbólicas: As funções hiperbólicas;  As derivadas das funções hiperbólicas. 
\\ 
 &  As técnicas de Integração: Integração por partes;  Integração por substituição trigonométrica;  Integração de potência das funções trigonométricas;  Integração por frações parciais. 
\\ 
 &  Integrais Impróprias: Formas indeterminadas;  A regra de L’Hôpital;  Integrais Impróprias. 
\\ 
 &  Sistema Polar: O sistema polar;  Gráficos em coordenadas polares;  Principais curvas polares. }

\def\PUDmetodo{Aulas expositórias que podem ser teóricas e/ou práticas, onde as práticas no laboratório serão marcadas no decorrer da disciplina.}

\def\PUDavalia{A avaliação será feita com aplicação de provas teóricas e/ou práticas, além da possibilidade de inclusão de trabalhos e seminários no decorrer da disciplina.}

\def\PUDbibbasica{ & \href{www.biblioteca.ifce.edu.br/index.asp?codigo_sophia=23890}{Thomas}, George B.;  Weir, Maurice D.;  Hass, Joel;  Giordano, Frank R.  Cálculo.  volume 1.  11º ed.  Editora Addison Wesley.  São Paulo, 2009.  ISBN: 9788588639317. 
\\ 
 & \href{www.biblioteca.ifce.edu.br/index.asp?codigo_sophia=24328}{Leithold}, Louis.  O cálculo com geometria analítica.  3º ed.  São Paulo: Harbra, 1994.  v.1.  ISBN: 8529400941. 
\\ 
 & \href{www.biblioteca.ifce.edu.br/index.asp?codigo_sophia=24186}{Leithold}, Louis.  O cálculo com geometria analítica.  3º ed.  São Paulo: Harbra, 1994.  v.2.  ISBN: 8529402065. }

\def\PUDbibcomplementar{ &  \href{www.biblioteca.ifce.edu.br/index.asp?codigo_sophia=23737}{Guidorizzi}, Hamilton L.  Um curso de cálculo.  Volume 1.  5º ed.  Rio de Janeiro, RJ : LTC, 2001.  635p.  ISBN: 9788521612599. 
\\ 
 &  \href{www.biblioteca.ifce.edu.br/index.asp?codigo_sophia=24164}{Guidorizzi}, Hamilton L.  Um curso de cálculo.  Volume 2.  5º ed.  Rio de Janeiro, RJ : LTC, 2001.  ISBN: 9788521612803. 
\\ 
 &  \href{www.biblioteca.ifce.edu.br/index.asp?codigo_sophia=23696}{Flemming}, Diva M.  Calculo A: funções, limite, derivadas e integração.  Volume Único.  6º ed.  São Paulo, SP : Pearson Prentice Hall, 2007.  449p.  ISBN: 857605115-X.
\\
 &  \href{www.biblioteca.ifce.edu.br/index.asp?codigo_sophia=62607}{Anton}, Howard.  Cálculo 2.  10º ed.  Porto Alegre, RS: Bookman, 2014.  669p.  ISBN: 9788582602454.
\\
 &  \href{www.biblioteca.ifce.edu.br/index.asp?codigo_sophia=62611}{Hughes-Hallett}, Deborah.  Cálculo aplicado.  4º ed.  Rio de Janeiro, RJ: LTC, 2012.  483p.  ISBN: 9788521620518. }

\def\PUDautor{Samuel Vieira Dias}

\def\PUDdata{2013-04-23}

\def\PUDrevisao{1}

\def\PUDrevisor{Samuel Vieira Dias}

\def\PUDdatarevisao{2017-05-09}

\PUDcorpo
